\section{Выпуск сертификата}\label{ACME.Order}

\subsection{Обработка заявки}\label{ACME.Order.Apply}

Клиент инициирует выпуск сертификата, отправляя 
\pdfexample{запрос}{acme/order1.txt} на узел \code{newOrder}. Тело запроса 
представляет собой подписанную заявку на выпуск сертификата 
(см.~\ref{ACME.Res.Order}) с обязательными полем \sstr{identifiers} и 
необязательными полями \sstr{notBefore} и \sstr{notAfter}. 
%
Поля заполняются по правилам, установленным в~\ref{ACME.Res.Order}.

\if0
\begin{codeblock}
POST /acme/new-order HTTP/1.1
Host: example.com
Content-Type: application/jose+json

{
  "protected": base64url({
    "alg": "BIGNS128",
    "kid": "https://example.com/acme/acct/evOfKhNU60wg",
    "nonce": "5XJ1L3lEkMG7tR6pA00clA",
    "url": "https://example.com/acme/new-order"
  }),
  "payload": base64url({
    "identifiers": [
      { "type": "dns", "value": "www.example.org" },
      { "type": "dns", "value": "example.org" }
    ],
    "notBefore": "2026-01-01T00:04:00+04:00",
    "notAfter": "2026-01-08T00:04:00+04:00"
  }),
  "signature": "H6ZXtGjTZyUnPeKn...wEA4TklBdh3e454g"
}
\end{codeblock}
\fi

Сервер должен возвратить ошибку, если он не может выполнить заявку в указанном 
виде. Сервер не должен выдавать сертификат с содержимым, не соответствующим 
заявке.
%
Если сервер требует изменить заявку определенным образом, он должен описать 
необходимые изменения, указывая в ответе соответствующий тип ошибки и прилагая 
подробности проблемы. 

Если сервер готов выдать запрошенный сертификат, то он дает
\pdfexample{ответ}{acme/order2.txt} с кодом статуса 201 (создано). Тело этого
ответа представляет собой объект заявки, в котором отражается запрос клиента и
перечисляются все авторизации, необходимые для выпуска сертификата.
%
Сервер возвращает URL заявки в заголовке \code{Location} ответа.

% use: https://errata.rfc-editor.org/eid5979/
   
\if0
\begin{codeblock}
HTTP/1.1 201 Created
Replay-Nonce: MYAuvOpaoIiywTezizk5vw
Link: <https://example.com/acme/directory>;rel="index"
Location: https://example.com/acme/order/TOlocE8rfgo

{
  "status": "pending",
  "expires": "2026-01-05T14:09:07.99Z",

  "notBefore": "2026-01-01T00:00:00Z",
  "notAfter": "2026-01-08T00:00:00Z",

  "identifiers": [
    { "type": "dns", "value": "www.example.org" },
    { "type": "dns", "value": "example.org" }
  ],

  "authorizations": [
    "https://example.com/acme/authz/PAniVnsZcis",
    "https://example.com/acme/authz/r4HqLzrSrpI"
  ],

  "finalize": "https://example.com/acme/order/TOlocE8rfgo/finalize"
} 
\end{codeblock}
\fi

Возвращая заявку, сервер информирует клиента, что если тот выполнит его
требования до истечения срока действия заявки (поле \sstr{expires}), 
то сервер обработает заявку и выдаст запрошенный сертификат. 
%
Требования сервера состоят в выполнении авторизаций, которые перечислены
в поле \sstr{authorizations} и имеют статус \sstr{pending}.
%
Если клиент не выполняет требуемые действия до истечения срока действия заявки, 
то сервер должен изменить ее статус на \sstr{invalid} и может удалить заявку.
%
Клиенты не должны делать предположений о порядке сортировки элементов 
массивов \sstr{identifiers} и \sstr{authorizations} в возвращаемой сервером
заявке.  

Как только клиент выполнил требования сервера, он должен отправить 
\pdfexample{POST-запрос}{acme/order3.txt} на URL-адрес финализации заявки. Тело 
запроса должно включать подписанный объект JSON с единственным полем: 

\begin{itemize}
\item
\sstr{csr} (обязательное, строка)~--- запрос на выпуск сертификата, 
сформированный в соответствии с~\ref{FMT.CSR}. DER-код запроса должен
быть дополнительно закодирован по правилам base64url.
\end{itemize}

\if0
\begin{codeblock}
POST /acme/order/TOlocE8rfgo/finalize HTTP/1.1
Host: example.com
Content-Type: application/jose+json

{
  "protected": base64url({
    "alg": "BIGNS128",
    "kid": "https://example.com/acme/acct/evOfKhNU60wg",
    "nonce": "MSF2j2nawWHPxxkE3ZJtKQ",
    "url": "https://example.com/acme/order/TOlocE8rfgo/finalize"
  }),
  "payload": base64url({
    "csr": "MIIBPTCBxAIBADBFMQ...FS6aKdZeGsysoCo4H9P",
  }),
  "signature": "uOrUfIIk5RyQ...nw62Ay1cl6AB"
}
\end{codeblock}
\fi

Запрос на выпуск сертификата определяет содержание будущего сертификата.
Запрос должен содержать тот же набор идентификаторов, что и исходная заявка
на узел \code{newOrder}.
%
Идентификаторы типа \sstr{dns} должны быть указаны в атрибуте
\code{extensionRequest}  в части, которая касается расширения 
\code{SubjectAltName}. Идентификаторы могут указываться в любом порядке.
%
Спецификации, которые определяют новые типы идентификаторов, должны указывать, 
в какой части запроса на выпуск сертификата эти идентификаторы должны 
появляться.

Запрос на финализацию заявки приведет к ошибке, если УЦ откажется 
выпустить сертификат по присланному запросу. Возможные причины отказа:
%
\begin{itemize}
\item
идентификаторы заявки и запроса различаются;
\item
авторизация не завершена;
\item
УЦ не готов включить в сертификат расширения, перечисленные в запросе.
\end{itemize}

В случае отказа сервер должен указать в своем ответе тип ошибки \sstr{badCSR},
а в поле \sstr{detail} подробностей проблемы конкретные причины отклонения 
запроса. После отказа сервер должен оставить заявку в состоянии \sstr{ready}, 
чтобы клиент мог отправить на адрес финализации уточненный запрос на выпуск 
сертификата.

Запрос финализации заявки приведет к ошибке, если заявка не находится в 
состоянии \sstr{ready}. В таких случаях сервер должен вернуть код статуса 
403 (запрещено) с типом ошибки \sstr{orderNotReady}. Затем клиент должен 
отправить запрос POST-as-GET на адрес заявки, чтобы получить ее текущее 
состояние. Статус заявки укажет клиенту, какие действия следует предпринять 
(см. ниже). 

Если запрос финализации заявки обработан успешно, то сервер дает 
\pdfexample{ответ}{acme/order4.txt} с кодом статуса 200 и включает в ответ 
обновленную заявку. Статус заявки указывает клиенту, какие действия следует 
предпринять:

\begin{itemize}
\item
\sstr{invalid}~--- сертификат не будет выдан. Заявку следует считать 
отмененной;

\item
\sstr{pending}~--- сервер считает, что клиент выполнил не все требования. 
Следует проверить массив \sstr{authorizations} на предмет наличия записей,  
которые находятся в состоянии ожидания;

\item
\sstr{ready}~--- сервер подтверждает выполнение требований и ожидает 
финализации. Следует отправить запрос на адрес финализации;
 
\item
\sstr{processing}~--- сертификат в процессе выпуска. Cледует отправить запрос 
POST-as-GET по истечении времени, указанного в заголовке \code{Retry-After}
ответа, если таковой имеется;

\item
\sstr{valid}~--- сервер выдал сертификат и указал его URL-адрес в поле
\sstr{certificate} заявки. Следует загрузить сертификат.
\end{itemize}

\if0
\begin{codeblock}
HTTP/1.1 200 OK
Replay-Nonce: CGf81JWBsq8QyIgPCi9Q9X
Link: <https://example.com/acme/directory>;rel="index"
Location: https://example.com/acme/order/TOlocE8rfgo

{
  "status": "valid",
  "expires": "2026-01-20T14:09:07.99Z",

  "notBefore": "2026-01-01T00:00:00Z",
  "notAfter": "2026-01-08T00:00:00Z",

  "identifiers": [
    { "type": "dns", "value": "www.example.org" },
    { "type": "dns", "value": "example.org" }
  ],

  "authorizations": [
    "https://example.com/acme/authz/PAniVnsZcis",
    "https://example.com/acme/authz/r4HqLzrSrpI"
  ],

  "finalize": "https://example.com/acme/order/TOlocE8rfgo/finalize",

  "certificate": "https://example.com/acme/cert/mAt3xBGaobw"
}
\end{codeblock}
\fi

\subsection{Предварительная авторизация}\label{ACME.Order.Preauthz}

Определенный выше процесс обработки заявки на выпуск сертификата предполагает, 
что объекты авторизации создаются реактивно, т.~е. в ответ на заявку. 
Некоторые серверы могут дать клиентам возможность выполнять авторизацию 
проактивно, т.~е. вне контекста конкретной заявки. Например, клиенту, который 
разворачивает виртуальные серверы для набора доменных имен, может быть удобно 
организовать авторизацию до запуска виртуальных серверов, а в момент запуска
всего лишь получить сертификаты. 

В некоторых случаях УЦ, использующий сервер ACME, может реализовать полностью 
внешний, не связанный с ACME процесс авторизации клиента. В таких случаях УЦ 
следует предоставить своему серверу ACME объекты авторизации, которые 
засчитываются как действительные при обработке любой заявки клиента.. 

Если УЦ разрешает предварительную авторизацию, то сервер ACME может добавить
в каталог ресурсов поле \sstr{newAuthz} с URL-адресом узла \code{newAuthz},
через который клиент может запрашивать предварительную авторизацию.

Клиент отправляет на узел \code{newAuthz} \pdfexample{запрос}{acme/order5.txt}, 
в теле которого содержится подписанный объект JSON с единственным полем:

\begin{itemize}
\item
\sstr{identifier} (обязательное, объект)~---
идентификатор, относительно которого запрашивается авторизация. 
%
Идентификатор представляет собой объект JSON с двумя обязательными 
полями:
\sstr{type} (строка)~--- тип идентификатора,
\sstr{value} (строка)~--- собственно идентификатор. 
\end{itemize}

\if0
POST /acme/new-authz HTTP/1.1
Host: example.com
Content-Type: application/jose+json

\begin{codeblock}
{
  "protected": base64url({
    "alg": "BIGNS128",
    "kid": "https://example.com/acme/acct/evOfKhNU60wg",
    "nonce": "uQpSjlRb4vQVCjVYAyyUWg",
    "url": "https://example.com/acme/new-authz"
  }),
  "payload": base64url({
    "identifier": {
      "type": "dns",
      "value": "example.org"
    }
  }),
  "signature": "nuSDISbWG8mMgE7H...QyVUL68yzf3Zawps"
}
\end{codeblock}
\fi

Идентификатор в запросе предварительной авторизации точно повторяется в объекте 
авторизации. Поэтому предварительная авторизация не может быть использована при 
выпуске сертификатов, содержащих подстановочные доменные имена.

Перед обработкой запроса авторизации серверу следует принять решение о 
допустимости выпуска сертификатов для данного идентификатора. Например, сервер 
может проверить тип идентификатора или осутствие идентификатора в определенном 
стоп-листе.
%
Если сервер отказывается выпускать сертификаты для данного идентификатора, то 
он должен возвратить ответ с кодом статуса 403 (запрещено) и подробностями
проблемы, описывающими причину отказа. 

Если сервер готов продолжить, он создает объект авторизации на основе данных, 
предоставленных клиентом. Объект должен содержать поля \sstr{identifier}, 
\sstr{status} и \sstr{challenges}, описанные в~\ref{ACME.Res.Authz}.
%
Поле \sstr{status} должно принимать значение \sstr{pending}, только если 
сервер не имеет дополнительной информации о статусе авторизации клиента.

Сервер выбирает новый URL-адрес авторизации и возвращает ответ с кодом 
статуса 201 (создано), выбранным адресом в заголовке \code{Location} 
и объектом авторизации в теле ответа. 
%
Затем клиент завершает авторизацию так, как это описано в~\ref{ACME.Authz}.

% skip: https://errata.rfc-editor.org/eid5861/:
% If a server receives a newAuthz request for an identifier where the 
% authorization object already exists, whether created by CA provisioning on 
% the ACME server or by the ACME server handling a previous newAuthz request 
% from a client, the server returns a 200 (OK) response with the existing 
% authorization URL in the Location header field and the existing JSON 
% authorization object in the body.

\subsection{Загрузка сертификата}\label{ACME.Order.Cert}
 
Чтобы загрузить выпущенный сертификат, клиент отправляет 
\pdfexample{запрос}{acme/order6.txt} POST-as-GET на URL-адрес сертификата.

В теле своего ответа сервер возвращает выпущенный сертификат,
а также сопровождающую его цепочку сертификатов. Каждый 
следующий сертификат цепочки является родительским для предыдущего.
Последний сертификат цепочки, точка доверия, может опускаться.     
%
Сертификаты задаются в формате PEM~\cite{RFC7468}.
%
PEM-код каждого сертификата начинается с новой строки.

% skip: The server MAY provide one or more link relation header fields
%  [RFC8288] with relation "alternate".  Each such field SHOULD express
%  an alternative certificate chain starting with the same end-entity
%  certificate.  This can be used to express paths to various trust
%  anchors.  Clients can fetch these alternates and use their own
%  heuristics to decide which is optimal.

\if0
\begin{codeblock}
POST /acme/cert/mAt3xBGaobw HTTP/1.1
Host: example.com
Content-Type: application/jose+json
Accept: application/pem-certificate-chain

{
  "protected": base64url({
    "alg": "BIGNS128",
    "kid": "https://example.com/acme/acct/evOfKhNU60wg",
    "nonce": "uQpSjlRb4vQVCjVYAyyUWg",
    "url": "https://example.com/acme/cert/mAt3xBGaobw"
  }),
  "payload": "",
  "signature": "nuSDISbWG8mMgE7H...QyVUL68yzf3Zawps"
}

HTTP/1.1 200 OK
Content-Type: application/pem-certificate-chain
Link: <https://example.com/acme/directory>;rel="index"

-----BEGIN CERTIFICATE-----
[End-entity certificate contents]
-----END CERTIFICATE-----
-----BEGIN CERTIFICATE-----
[Issuer certificate contents]
-----END CERTIFICATE-----
-----BEGIN CERTIFICATE-----
[Other certificate contents]
-----END CERTIFICATE-----
\end{codeblock}
\fi

Выпущенный сертификат является неизменяемым ресурсом. Если клиенту
нужно обновить сертификат, он должен запустить процесс выпуска заново.
%
Поскольку сертификат нельзя изменить после выпуска, сервер 
может включить кэширование ресурса, добавив заголовки \code{Expires} и 
\code{Cache-Control} и указав в них момент времени в отдаленном будущем. 
Время в заголовках не имеет отношения к сроку действия сертификата. 

% skip: The ACME client MAY request other formats by including an Accept
%  header field [RFC7231] in its request.  For example, the client could
%  use the media type "application/pkix-cert" [RFC2585] or "application/
%  pkcs7-mime" [RFC5751] to request the end-entity certificate in DER
%  format.  Server support for alternate formats is OPTIONAL.  For
%  formats that can only express a single certificate, the server SHOULD
%  provide one or more "Link: rel="up"" header fields pointing to an
%  issuer or issuers so that ACME clients can build a certificate chain
%  as defined in TLS (see Section 4.4.2 of [RFC8446]).

